% ---
% titre: Problème d'échelle - la carte de la Suisse
% theme: FA
% chapitre: Proportionnalité
% sous_chapitre: Résoudre des problèmes
% niveau: [10e]
% type: EVAL
% tags: [proportionnalite, p-question-TS]
% difficulte: 1
% corrige: true
% ---

\question
Voici une carte de la Suisse à l'échelle 1:2\,400\,000. 

\begin{center}
\includegraphics[width=\linewidth]{IMG_FA-Proportionnalite-CarteSuisse.pdf}
\end{center}

\begin{parts}

\vspace{20pt}\part[2\half]
Calcule la distance réelle, à vol d'oiseau, entre les villes de Lausanne et Zürich.
	
\vspace{10pt}{\footnotesize\raggedleft Espace pour ta démarche et tes calculs\par}
\begin{solutionorgrid}[140pt]

	\vspace{5pt}
	distance sur la carte: 7,1\,cm ($\pm$ 2\,mm) \hfill\textbf{(0,5\,pt)}

	\vspace{15pt}
	\def\arraystretch{1.5}
	\begin{tabularx}{0.6\textwidth}{l|C|C|}
	Distance sur la carte [cm]&1&7,1\\
	\hline
	Distance réelle [cm]&2\,400\,000&\\
	\end{tabularx}
	\def\arraystretch{1}
	
	\hfill\textbf{(0,5pt)} pour la construction du tableau

	\vspace{15pt}
	$7,1 \cdot 2\,400\,000 : 1 = 17\,040\,000$\,cm  \hfill calcul \textbf{(0,5pt)} et réponse \textbf{(0,5pt)}
	\end{solutionorgrid}
	
\begin{tcolorbox}[enhanced jigsaw, interior hidden, colframe=box, parbox=false]%
Réponse: 
\sol{\vspace{20pt}}{Zürich se trouve à 17\,040\,000\,cm (170.4\,km) de Lausanne. \textbf{(0,5pt)}}
\end{tcolorbox}
	
\newpage
\part[8] Le train intercity 1 relie les villes de Lausanne, Fribourg, Bern et Zürich. Il met 50\,min pour relier Lausanne et Fribourg. Combien de temps dure le voyage Lausanne-Zürich considérant que le train roule à vitesse moyenne constante pendant tout le trajet?
	
\vspace{10pt}
{\small Indice: commence par calculer la distance entre Lausanne et Fribourg avant de calculer le temps de trajet entre Lausanne et Zürich.}

\vspace{10pt}{\footnotesize\raggedleft Espace pour ta démarche et tes calculs\par}
\begin{solutionorgrid}[300pt]
	
	\vspace{5pt}
	Distance Lausanne Fribourg: 
	
	\vspace{5pt}distance sur la carte: 2,1\,cm ($\pm$ 2\,mm) \hfill\textbf{(0,5\,pt)}
	
	\vspace{5pt}
	\def\arraystretch{1.5}
	\begin{tabularx}{0.6\textwidth}{l|C|C|}
	Distance sur la carte [cm]& 1 & 2.1 \\
	\hline
	Distance réelle [cm]& 2\,400\,000 & \\
	\end{tabularx}
	
	\hfill\textbf{(0,5pt)} pour la construction du tableau
	
	\vspace{5pt}
	$2,1 \cdot 2\,400\,000 :1 = 5\,040\,000$\,cm \hfill calcul \textbf{(0,5pt)} et réponse \textbf{(0,5pt)}

	
	\vspace{15pt}
	Temps du trajet:
	
	\vspace{5pt}
	\begin{tabularx}{0.6\textwidth}{l|C|C|}
	Durée [min]& 50 & \\
	\hline
	Distance [cm] & 5\,040\,000 & 17\,040\,000 \\
	\end{tabularx}
	\def\arraystretch{1}
	
	\hfill\textbf{(0,5pt)} pour la construction du tableau
	
	\vspace{5pt}
	$50 \cdot 17\,040\,000 : 5\,040\,000 = 169$\,min \hfill calcul \textbf{(0,5pt)} et réponse \textbf{(0,5pt)}
	\end{solutionorgrid}
	
\begin{tcolorbox}[enhanced jigsaw, interior hidden, colframe=box, parbox=false]%
Réponse: 
\sol{\vspace{20pt}}{Le trajet Lausanne-Zürich dure 169\,min (2h et 49 minutes).\textbf{(0,5pt)}}
\end{tcolorbox}
\end{parts}

